% WasmSentinel — ACM AsiaCCS 2027 (Cycle 2, deadline 2026-12-11 AoE)
% sigconf, double-blind. Body <= 12 pages, references + appendix <= 10 pages.
% Every number in this paper must trace to evaluation/results/2026-11/ (see CLAIMS.md).

\documentclass[sigconf,anonymous,review]{acmart}

\usepackage{booktabs}
\usepackage{subcaption}
\usepackage{listings}
\usepackage{xcolor}
\usepackage{tikz}
\usetikzlibrary{shapes,arrows,positioning,fit,backgrounds}
\usepackage{siunitx}
\sisetup{detect-all,group-separator={,},group-minimum-digits=4}

\lstdefinestyle{rust}{
  language=C, basicstyle=\ttfamily\footnotesize, keywordstyle=\color{blue}\bfseries,
  commentstyle=\color{green!60!black}, stringstyle=\color{orange},
  numbers=left, numberstyle=\tiny\color{gray}, breaklines=true, frame=single, captionpos=b}

% Review build: no copyright block, no ACM reference.
\setcopyright{none}
\settopmatter{printacmref=false}
\renewcommand\footnotetextcopyrightpermission[1]{}
\pagestyle{plain}
\acmConference[AsiaCCS '27]{ACM Asia Conference on Computer and Communications Security}{2027}{TBD}

\newcommand{\sys}{\textsc{WasmSentinel}\xspace}
\usepackage{xspace}
% Numbers live in tables/numbers.tex, generated from evaluation/results/2026-11 (make paper-numbers).
\IfFileExists{tables/numbers.tex}{% Generated by evaluation/scripts/paper_numbers.py from evaluation/results/2026-11
% commits: e1=ba7af70ee77e, e2=ba7af70ee77e, e3=ba7af70ee77e, e4=ba7af70ee77e, e5=ba7af70ee77e
% Do not edit; run `make numbers` in paper-asiaccs/.
\newcommand{\isoContained}{11}
\newcommand{\isoTotal}{11}
\newcommand{\isoModules}{10}
\newcommand{\isoMemoryCapMiB}{64}
\newcommand{\isoFuel}{50,000,000}
\newcommand{\isoTimeout}{100ms}
\newcommand{\parityInputs}{1,400}
\newcommand{\parityMaxScoreDiff}{2.0e-06}
\newcommand{\provCpu}{AMD Ryzen 9 5950X}
\newcommand{\provKernel}{7.0.0-34-generic}
\newcommand{\eOneNormal}{5,205}
\newcommand{\eOneAttack}{746}
\newcommand{\eOneFolds}{5}
\newcommand{\eOneSeeds}{3}
\newcommand{\eOneAuprcMlp}{0.730}
\newcommand{\eOneAuprcStdMlp}{0.056}
\newcommand{\eOneAurocMlp}{0.946}
\newcommand{\eOneRecallFprFiveMlp}{0.718}
\newcommand{\eOneRecallFprOneMlp}{0.345}
\newcommand{\eOneFOneOpFiveMlp}{0.691}
\newcommand{\eOneAuprcStide}{0.437}
\newcommand{\eOneAuprcStdStide}{0.043}
\newcommand{\eOneAurocStide}{0.787}
\newcommand{\eOneRecallFprFiveStide}{0.377}
\newcommand{\eOneRecallFprOneStide}{0.157}
\newcommand{\eOneFOneOpFiveStide}{0.344}
\newcommand{\eOneAuprcIForest}{0.218}
\newcommand{\eOneAuprcStdIForest}{0.007}
\newcommand{\eOneAurocIForest}{0.745}
\newcommand{\eOneRecallFprFiveIForest}{0.046}
\newcommand{\eOneRecallFprOneIForest}{0.002}
\newcommand{\eOneFOneOpFiveIForest}{0.060}
\newcommand{\eOneAuprcRf}{0.897}
\newcommand{\eOneAuprcStdRf}{0.028}
\newcommand{\eOneAurocRf}{0.982}
\newcommand{\eOneRecallFprFiveRf}{0.908}
\newcommand{\eOneRecallFprOneRf}{0.700}
\newcommand{\eOneFOneOpFiveRf}{0.636}
\newcommand{\eOneAuprcLstm}{0.825}
\newcommand{\eOneAuprcStdLstm}{0.061}
\newcommand{\eOneAurocLstm}{0.965}
\newcommand{\eOneRecallFprFiveLstm}{0.858}
\newcommand{\eOneRecallFprOneLstm}{0.488}
\newcommand{\eOneFOneOpFiveLstm}{0.760}
\newcommand{\eOneDiffStide}{+0.293}
\newcommand{\eOneDiffCiLoStide}{+0.257}
\newcommand{\eOneDiffCiHiStide}{+0.328}
\newcommand{\eOneDiffIForest}{+0.512}
\newcommand{\eOneDiffCiLoIForest}{+0.480}
\newcommand{\eOneDiffCiHiIForest}{+0.544}
\newcommand{\eOneDiffRf}{-0.167}
\newcommand{\eOneDiffCiLoRf}{-0.199}
\newcommand{\eOneDiffCiHiRf}{-0.135}
\newcommand{\eOneDiffLstm}{-0.095}
\newcommand{\eOneDiffCiLoLstm}{-0.126}
\newcommand{\eOneDiffCiHiLstm}{-0.064}
\newcommand{\eOneWasmNumpyMaxDiff}{1.1e-06}
\newcommand{\eOneFamAdduserMlp}{0.79}
\newcommand{\eOneFamHydraFtpMlp}{0.80}
\newcommand{\eOneFamHydraSshMlp}{0.87}
\newcommand{\eOneFamJavaMeterpreterMlp}{0.80}
\newcommand{\eOneFamMeterpreterMlp}{0.71}
\newcommand{\eOneFamWebShellMlp}{0.66}
\newcommand{\eOneFamAdduserStide}{0.78}
\newcommand{\eOneFamHydraFtpStide}{0.85}
\newcommand{\eOneFamHydraSshStide}{0.84}
\newcommand{\eOneFamJavaMeterpreterStide}{0.85}
\newcommand{\eOneFamMeterpreterStide}{0.68}
\newcommand{\eOneFamWebShellStide}{0.78}
\newcommand{\eOneFamAdduserIForest}{0.04}
\newcommand{\eOneFamHydraFtpIForest}{0.00}
\newcommand{\eOneFamHydraSshIForest}{0.03}
\newcommand{\eOneFamJavaMeterpreterIForest}{0.05}
\newcommand{\eOneFamMeterpreterIForest}{0.13}
\newcommand{\eOneFamWebShellIForest}{0.06}
\newcommand{\eOneFamAdduserRf}{0.96}
\newcommand{\eOneFamHydraFtpRf}{0.99}
\newcommand{\eOneFamHydraSshRf}{0.99}
\newcommand{\eOneFamJavaMeterpreterRf}{0.99}
\newcommand{\eOneFamMeterpreterRf}{0.97}
\newcommand{\eOneFamWebShellRf}{0.95}
\newcommand{\eOneFamAdduserLstm}{0.83}
\newcommand{\eOneFamHydraFtpLstm}{0.91}
\newcommand{\eOneFamHydraSshLstm}{0.92}
\newcommand{\eOneFamJavaMeterpreterLstm}{0.91}
\newcommand{\eOneFamMeterpreterLstm}{0.87}
\newcommand{\eOneFamWebShellLstm}{0.76}
\newcommand{\eTwoTraces}{5,951}
\newcommand{\eTwoMaxDiffNumpyWasm}{1.9e-06}
\newcommand{\eTwoMaxDiffNativeWasm}{5.0e-09}
\newcommand{\eTwoTraceLen}{343}
\newcommand{\eTwoTracesTimed}{662}
\newcommand{\eTwoUsPerCallNumpy}{164.6}
\newcommand{\eTwoUsPerCallNative}{26.0}
\newcommand{\eTwoUsPerCallWasm}{126.3}
\newcommand{\eTwoEngineOverheadUs}{2.1}
\newcommand{\eTwoMlJsonUs}{10.2}
\newcommand{\eTwoMlRawUs}{4.6}
\newcommand{\eThreeAchievedJsonOneK}{964}
\newcommand{\eThreePFiftyJsonOneK}{118}
\newcommand{\eThreePNinetyNineJsonOneK}{183}
\newcommand{\eThreePNineNineNineJsonOneK}{325}
\newcommand{\eThreeRssMiBJsonOneK}{90}
\newcommand{\eThreeDetectPFiftyJsonOneK}{105}
\newcommand{\eThreeDetectPNinetyNineJsonOneK}{149}
\newcommand{\eThreeAchievedBinOneK}{961}
\newcommand{\eThreePFiftyBinOneK}{101}
\newcommand{\eThreePNinetyNineBinOneK}{144}
\newcommand{\eThreePNineNineNineBinOneK}{240}
\newcommand{\eThreeRssMiBBinOneK}{92}
\newcommand{\eThreeDetectPFiftyBinOneK}{89}
\newcommand{\eThreeDetectPNinetyNineBinOneK}{119}
\newcommand{\eThreeAchievedJsonTenK}{9,689}
\newcommand{\eThreePFiftyJsonTenK}{378}
\newcommand{\eThreePNinetyNineJsonTenK}{1,853}
\newcommand{\eThreePNineNineNineJsonTenK}{4,686}
\newcommand{\eThreeRssMiBJsonTenK}{99}
\newcommand{\eThreeDetectPFiftyJsonTenK}{81}
\newcommand{\eThreeDetectPNinetyNineJsonTenK}{131}
\newcommand{\eThreeAchievedBinTenK}{9,688}
\newcommand{\eThreePFiftyBinTenK}{304}
\newcommand{\eThreePNinetyNineBinTenK}{1,246}
\newcommand{\eThreePNineNineNineBinTenK}{3,585}
\newcommand{\eThreeRssMiBBinTenK}{100}
\newcommand{\eThreeDetectPFiftyBinTenK}{70}
\newcommand{\eThreeDetectPNinetyNineBinTenK}{112}
\newcommand{\eThreeAchievedJsonFiftyK}{14,996}
\newcommand{\eThreePFiftyJsonFiftyK}{334,813}
\newcommand{\eThreePNinetyNineJsonFiftyK}{425,654}
\newcommand{\eThreePNineNineNineJsonFiftyK}{445,640}
\newcommand{\eThreeRssMiBJsonFiftyK}{178}
\newcommand{\eThreeDetectPFiftyJsonFiftyK}{82}
\newcommand{\eThreeDetectPNinetyNineJsonFiftyK}{117}
\newcommand{\eThreeAchievedBinFiftyK}{18,466}
\newcommand{\eThreePFiftyBinFiftyK}{270,773}
\newcommand{\eThreePNinetyNineBinFiftyK}{352,795}
\newcommand{\eThreePNineNineNineBinFiftyK}{366,660}
\newcommand{\eThreeRssMiBBinFiftyK}{189}
\newcommand{\eThreeDetectPFiftyBinFiftyK}{71}
\newcommand{\eThreeDetectPNinetyNineBinFiftyK}{95}
\newcommand{\eThreeBinGainPct}{23}
\newcommand{\eFourTputNginxNone}{178,007}
\newcommand{\eFourCpuNginxNone}{0.0}
\newcommand{\eFourRssNginxNone}{0}
\newcommand{\eFourTputNginxWs}{170,071}
\newcommand{\eFourDeltaNginxWs}{-4.4}
\newcommand{\eFourDeltaCiLoNginxWs}{-7.2}
\newcommand{\eFourDeltaCiHiNginxWs}{-1.7}
\newcommand{\eFourCpuNginxWs}{18.4}
\newcommand{\eFourRssNginxWs}{111}
\newcommand{\eFourTputNginxFalco}{161,597}
\newcommand{\eFourDeltaNginxFalco}{-9.2}
\newcommand{\eFourDeltaCiLoNginxFalco}{-11.8}
\newcommand{\eFourDeltaCiHiNginxFalco}{-6.6}
\newcommand{\eFourCpuNginxFalco}{82.0}
\newcommand{\eFourRssNginxFalco}{368}
\newcommand{\eFourTputNginxTetragon}{148,390}
\newcommand{\eFourDeltaNginxTetragon}{-16.6}
\newcommand{\eFourDeltaCiLoNginxTetragon}{-18.1}
\newcommand{\eFourDeltaCiHiNginxTetragon}{-15.2}
\newcommand{\eFourCpuNginxTetragon}{282.1}
\newcommand{\eFourRssNginxTetragon}{195}
\newcommand{\eFourTputRedisNone}{960,059}
\newcommand{\eFourCpuRedisNone}{0.0}
\newcommand{\eFourRssRedisNone}{0}
\newcommand{\eFourTputRedisWs}{927,451}
\newcommand{\eFourDeltaRedisWs}{-3.4}
\newcommand{\eFourDeltaCiLoRedisWs}{-5.0}
\newcommand{\eFourDeltaCiHiRedisWs}{-1.8}
\newcommand{\eFourCpuRedisWs}{8.1}
\newcommand{\eFourRssRedisWs}{97}
\newcommand{\eFourTputRedisFalco}{919,601}
\newcommand{\eFourDeltaRedisFalco}{-4.2}
\newcommand{\eFourDeltaCiLoRedisFalco}{-7.2}
\newcommand{\eFourDeltaCiHiRedisFalco}{-1.2}
\newcommand{\eFourCpuRedisFalco}{7.4}
\newcommand{\eFourRssRedisFalco}{369}
\newcommand{\eFourTputRedisTetragon}{945,367}
\newcommand{\eFourDeltaRedisTetragon}{-1.5}
\newcommand{\eFourDeltaCiLoRedisTetragon}{-3.5}
\newcommand{\eFourDeltaCiHiRedisTetragon}{+0.5}
\newcommand{\eFourCpuRedisTetragon}{11.9}
\newcommand{\eFourRssRedisTetragon}{96}
\newcommand{\eFourTputPgNone}{26,450}
\newcommand{\eFourCpuPgNone}{0.0}
\newcommand{\eFourRssPgNone}{0}
\newcommand{\eFourTputPgWs}{23,799}
\newcommand{\eFourDeltaPgWs}{-10.0}
\newcommand{\eFourDeltaCiLoPgWs}{-14.4}
\newcommand{\eFourDeltaCiHiPgWs}{-5.5}
\newcommand{\eFourCpuPgWs}{149.6}
\newcommand{\eFourRssPgWs}{161}
\newcommand{\eFourTputPgFalco}{23,712}
\newcommand{\eFourDeltaPgFalco}{-10.2}
\newcommand{\eFourDeltaCiLoPgFalco}{-15.3}
\newcommand{\eFourDeltaCiHiPgFalco}{-5.1}
\newcommand{\eFourCpuPgFalco}{80.7}
\newcommand{\eFourRssPgFalco}{365}
\newcommand{\eFourTputPgTetragon}{25,907}
\newcommand{\eFourDeltaPgTetragon}{-1.9}
\newcommand{\eFourDeltaCiLoPgTetragon}{-7.9}
\newcommand{\eFourDeltaCiHiPgTetragon}{+4.1}
\newcommand{\eFourCpuPgTetragon}{14.2}
\newcommand{\eFourRssPgTetragon}{104}
\newcommand{\eFiveTrials}{20}
\newcommand{\eFiveDetectedWsShadow}{20}
\newcommand{\eFiveLatencyWsShadow}{2.9}
\newcommand{\eFiveDetectedWsTmpExec}{20}
\newcommand{\eFiveLatencyWsTmpExec}{0.8}
\newcommand{\eFiveDetectedWsConnect}{20}
\newcommand{\eFiveLatencyWsConnect}{2.3}
\newcommand{\eFiveRssWs}{97}
\newcommand{\eFiveCpuWs}{7.1}
\newcommand{\eFiveDetectedFalcoShadow}{20}
\newcommand{\eFiveLatencyFalcoShadow}{2.8}
\newcommand{\eFiveDetectedFalcoTmpExec}{20}
\newcommand{\eFiveLatencyFalcoTmpExec}{0.7}
\newcommand{\eFiveDetectedFalcoConnect}{20}
\newcommand{\eFiveLatencyFalcoConnect}{1.8}
\newcommand{\eFiveRssFalco}{371}
\newcommand{\eFiveCpuFalco}{1.2}
\newcommand{\eFiveDetectedTetragonShadow}{20}
\newcommand{\eFiveLatencyTetragonShadow}{2.7}
\newcommand{\eFiveDetectedTetragonTmpExec}{20}
\newcommand{\eFiveLatencyTetragonTmpExec}{0.8}
\newcommand{\eFiveDetectedTetragonConnect}{20}
\newcommand{\eFiveLatencyTetragonConnect}{1.7}
\newcommand{\eFiveRssTetragon}{154}
\newcommand{\eFiveCpuTetragon}{0.8}
}{}

\begin{document}

\title{WasmSentinel: What It Costs and What It Buys to Run Host-Based Detection Rules in a WebAssembly Sandbox}

\author{Anonymous Authors}
\affiliation{\institution{Anonymous Institution}\country{Anonymous Country}}

\begin{abstract}
Host-based detection rules see everything a machine does and, in the
eBPF monitors deployed today, run as trusted code inside a privileged
agent. We ask what it costs, and what it buys, to run them as untrusted
WebAssembly modules instead. \sys collects process, file and network
events with a fixed set of eBPF programs and hands each event to every
rule inside its own Wasm instance under a memory cap, an instruction
budget and a per-call deadline. We measure the boundary rather than
assume it. A rule call costs \eTwoEngineOverheadUs{}~\si{\micro\second}
empty and \eTwoMlRawUs{}~\si{\micro\second} for a real syscall detector;
the pipeline sustains \eThreeAchievedBinFiftyK{} events per second; and
under nginx, redis and PostgreSQL the application cost
(\eFourDeltaNginxWs{}, \eFourDeltaRedisWs{}, \eFourDeltaPgWs{}\%) is in
the range of Falco's, with a smaller memory footprint and, for
syscall-heavy workloads, a larger CPU one whose cause we identify. Ten
adversarial modules that try to read files, open sockets, exhaust memory
or spin forever are contained in all \isoTotal{} attempts. Because the
deployed module reproduces its trained models to \eTwoMaxDiffNumpyWasm{}
(an MLP) and \eOneWasmSklearnMaxDiff{} (a random forest), we can report
its accuracy on ADFA-LD with every score computed inside the sandbox: an
AUPRC of \eOneAuprcRfWasm{} for the forest, above an LSTM and the MLP the
module first shipped with (\eOneAuprcMlp{}), at no higher per-event cost.
All numbers in the paper are generated from released result files.
\end{abstract}

\begin{CCSXML}
<ccs2012>
<concept><concept_id>10002978.10002991.10002992</concept_id><concept_desc>Security and privacy~Intrusion detection systems</concept_desc><concept_significance>500</concept_significance></concept>
<concept><concept_id>10011007.10011006.10011066</concept_id><concept_desc>Software and its engineering~Runtime environments</concept_desc><concept_significance>300</concept_significance></concept>
</ccs2012>
\end{CCSXML}
\ccsdesc[500]{Security and privacy~Intrusion detection systems}
\ccsdesc[300]{Software and its engineering~Runtime environments}
\keywords{host intrusion detection, eBPF, WebAssembly, sandboxing, syscall analysis}

\maketitle

\section{Introduction}
\label{sec:intro}

A host intrusion detection rule is a small program with a large view. It
sees every process a machine starts, every file it opens and every
connection it makes, and it decides which of them deserve an alert. In the
monitors that run on most Linux fleets today the rule is also
\emph{trusted}: Falco evaluates its rules and loads its plugins inside one
privileged agent~\cite{falco}, Tetragon compiles its policies into the
kernel and a privileged agent~\cite{tetragon}, and a mistake in a rule,
or a rule from a source that should not have been trusted, runs with the
agent's view and the agent's power. Rules come from many places---vendor
packs, community repositories, a security team's own detectors, and
increasingly a model trained by someone else---and they change far more
often than the agent does.

This paper asks a narrow question about that arrangement. Suppose the
rules were untrusted and ran inside a WebAssembly sandbox, each in its own
instance, with no access to the host beyond the event it is handed and
with a memory cap, an instruction budget and a wall-clock deadline on
every call. \emph{What does that cost, and what does it buy?} The
question is narrow on purpose. WebAssembly's isolation properties are
established~\cite{wasm-spec,rlbox,wasmboxc}; eBPF-based collection is
established~\cite{falco,tetragon,tracee}; what is missing is a measurement
of the boundary between them in this particular use, where the sandboxed
code is fed by the kernel at tens of thousands of events per second and
the alternative it is compared against is not ``no sandbox'' but two
production systems.

We built \sys to take that measurement. Six fixed eBPF programs produce a
fixed \SI{592}{B} record per event; a Go daemon queues the records and
hands each one to every rule through a five-function ABI; a rule is a Wasm
module that returns a severity. Two rules ship with it: a pattern rule of
the kind one writes in Falco's filter language, and a syscall-sequence
detector that combines a STIDE $n$-gram database~\cite{stide} with a
small MLP whose weights are embedded at build time. The design choices
that affect the measurement are few: one long-lived instance per rule,
limits attached to the instance's store, a trapped rule dropped and
re-instantiated, and a choice between a JSON encoding of the event and the
kernel record passed verbatim.

The evaluation makes three claims, each with the experiment that supports
it.

\begin{enumerate}
\item \textbf{Cost (E2--E5).} A sandboxed rule call costs
  \eTwoEngineOverheadUs{}~\si{\micro\second} empty and
  \eTwoMlRawUs{}~\si{\micro\second} for the real detector on a process
  event; end to end, the pipeline sustains \eThreeAchievedBinFiftyK{}
  syscall events per second with the binary encoding
  (\eThreeAchievedJsonFiftyK{} with JSON) at a median latency of
  \eThreePFiftyBinTenK{}~\si{\micro\second} at \num{10000}~events/s. Under
  nginx, redis and PostgreSQL, \sys costs the application
  \eFourDeltaNginxWs{}, \eFourDeltaRedisWs{} and \eFourDeltaPgWs{}\% of
  throughput, in the range of Falco's \eFourDeltaNginxFalco{},
  \eFourDeltaRedisFalco{} and \eFourDeltaPgFalco{}\%, with a smaller
  resident footprint and, under PostgreSQL, a larger CPU one; the three
  systems detect the same three actions in every trial with alert
  latencies within a millisecond of each other. The CPU gap has a named
  cause---every event reaches every rule, and the ML rule recomputes its
  features per event---and is a lever on the rule side, not a sandbox
  tax.
\item \textbf{Isolation (E6).} \isoModules{} modules written to read
  files, open sockets, leak the environment, flood output, exhaust
  memory, spin forever, overflow the stack or exit the process were all
  contained under the default limits in \isoContained{} of \isoTotal{}
  attempts, with the neighbouring rule unaffected and the next call
  succeeding. We report the one capability a rule does have (the clock)
  and the one case the runtime classifies less precisely than we would
  like.
\item \textbf{Equivalence (E1--E2).} The module the daemon loads produces,
  to within \eTwoMaxDiffNumpyWasm{}, the same score as the reference
  implementation that was trained, on all \eTwoTraces{} ADFA-LD traces.
  Because of that, the accuracy we report is the module's own: every test
  score in a \eOneFolds-fold, \eOneSeeds-seed evaluation is produced
  inside the sandbox. That accuracy is an AUPRC of \eOneAuprcMlp{},
  above STIDE and an Isolation Forest and \emph{below} a random forest
  (\eOneAuprcRf{}) and an LSTM (\eOneAuprcLstm{}) on the same data. We
  say so because the point of E1 is the protocol, and because nothing in
  the sandbox prefers the weaker model.
\end{enumerate}

We do not claim that \sys is the earliest system to combine eBPF and
WebAssembly, nor that its detector is state of the art, nor that its
event path is copy-free. The nearest prior work combines the two
technologies in other directions (\S\ref{sec:related}); the detector is
deliberately ordinary; and the binary encoding removes a parse, not a
copy. What the paper contributes is the measurement, taken on the deployed
artifact, with every number regenerated from the result files we release
(\S\ref{sec:openscience}).

\section{Background}
\label{sec:background}

\paragraph{eBPF as a sensor.}
{\sloppy eBPF lets a user-space program attach small verified programs to kernel
hooks---tracepoints such as \texttt{sched\_process\_exec}, kprobes on
functions such as \texttt{tcp\_connect}, and the raw system-call entry
point---and ship data to user space through a ring buffer~\cite{ebpf-intro}.
The verifier bounds what such a program may do (no unbounded loops, no
unchecked memory access), which is what makes it safe to run in the kernel
and also what makes it a poor place for a detector with state, a model, or
a learned database. Production monitors therefore split the work: eBPF
collects, user space decides~\cite{falco,tetragon,tracee}. \sys keeps that
split and changes only what ``user space decides'' runs in.\par}

\paragraph{WebAssembly as a sandbox.}
A WebAssembly module~\cite{wasm-spec} computes over a single linear memory
that it cannot address outside of, calls the host only through functions
the host chose to import, and has no notion of a file, socket or process
unless an import provides one. WASI~\cite{wasi} is the standard set of such
imports; a host may link all of it, part of it, or an implementation that
grants nothing. Runtimes such as Wasmtime~\cite{wasmtime} compile modules
ahead of time to native code and add two controls the specification does
not: \emph{fuel}, a per-store instruction budget that traps when exhausted,
and \emph{epochs}, a counter the host increments on a timer that traps a
running module when a deadline passes. Together with a resource limiter on
\texttt{memory.grow}, these give a host three knobs---memory, instructions,
wall clock---that it can set per instance and per call. The cost of the
boundary is a copy into linear memory and a native call; prior work has
pushed that transition toward zero~\cite{wasmboxc} and used the same
substrate to isolate browser libraries~\cite{rlbox} and proxy
plugins~\cite{proxy-wasm}. What has not been measured is what it costs
when the code behind the boundary is a host-intrusion-detection rule fed by
the kernel at tens of thousands of events per second; that is the gap
\S\ref{sec:eval} fills.

\paragraph{Syscall-sequence detection.}
The detector \sys ships is deliberately conventional. STIDE~\cite{stide}
records every $n$-gram of syscall numbers seen in normal traces and flags
an $n$-gram it has not seen; it is cheap, needs only normal data, and has
well-known blind spots. Supervised models on windowed features or on the raw
sequence~\cite{creech,kim-lstm} do better on ADFA-LD~\cite{adfa-ld}, the
benchmark we use, at a higher cost. We use STIDE as a cheap first stage
and a small MLP as the scorer, and compare the result to both families in
\S\ref{sec:eval:e1} without claiming it is the best choice.

\section{Threat Model}
% Adversaries: (A1) the monitored workload (attacker on the host) producing
% events; (A2) a malicious or buggy detection rule (third-party module).
% Trusted: kernel, eBPF programs, the host process. Out of scope: kernel
% exploits, side channels in wasmtime, the collector being disabled by root.
% State exactly which capabilities a rule has (clock) and lacks (fs, net, env,
% other rules' memory) -- matches E6's table.
\label{sec:threat}
TODO

\section{Design}
% Collector (6 hooks, record layout 592 B) -> pipeline (back-pressure queue)
% -> engine (one long-lived instance per rule; alloc/dealloc/detect ABI;
% memory cap, fuel, epoch deadline; trap => reinstantiate; zero-capability
% WASI) -> outputs. Event encodings: JSON vs binary record (sentinel-abi).
% Rules: suspicious_shell (keyword/port/path), ml_detector (STIDE fast path +
% MLP deep path on feature set v1; training path train_sequence).
\label{sec:design}
TODO

\section{Implementation}
\label{sec:impl}
% Line counts: `grep -v '^\s*(//|$)' | wc -l` at the commit named in
% tables/numbers.tex; generated files (bpf2go bindings, weights.rs) excluded
% where stated. Module sizes from `ls -l rules/wasm`.

The daemon is about 3,400 lines of Go (plus 1,600 of tests) on
wasmtime-go~v17~\cite{wasmtime} and cilium/ebpf; the kernel side is 190
lines of C compiled with clang and embedded through \texttt{bpf2go}, so the
binary carries its own eBPF objects and the record layout is checked
against the C struct at build time. Rules are Rust compiled to
\texttt{wasm32-wasip1}. A 620-line crate, \texttt{sentinel-abi}, holds the
record layout, the binary encoding and the \texttt{EventFields} trait and is
shared by every rule; \texttt{suspicious\_shell} is 330 lines and
\SI{91}{kB} of Wasm, \texttt{ml\_detector} is 1,900 lines plus a generated
19,000-line weights file and \SI{334}{kB} of Wasm. The ten adversarial
modules of \S\ref{sec:eval:e6} total 220 lines. Nothing in the host is
specific to Rust: any toolchain that can export the five functions of
\S\ref{sec:design:engine} produces a loadable rule.

Three implementation details matter for the measurements. First, the
engine reuses one guest buffer per rule across calls (it asks the guest to
allocate once and grows only when an event is larger), so the steady-state
cost of a call is one copy plus the invocation. Second, fuel and epoch
enforcement are enabled at engine construction; both add a small, fixed
per-instruction and per-tick cost that E2's trivial-guest benchmark
includes. Third, the feature extractor in \texttt{ml\_detector} is a
line-by-line port of the Python function used for training, including its
iteration order over hash maps, and a golden test with 18 recorded inputs
runs in \texttt{cargo test} so that the parity of \S\ref{sec:eval:e2} is
checked on every build. The same source tree builds the replay collector,
the metrics endpoint and the evaluation harnesses; every table in this
paper is regenerated by one \texttt{make} target from the JSON files it
cites.

\section{Evaluation}
% Setup: hardware/kernel from provenance blocks in the JSONs.
% E2 consistency + single-call cost (table) -> E3 end-to-end latency/throughput,
% JSON vs binary (figure) -> E4 application overhead vs Falco/Tetragon (table)
% -> E5 same detections, three systems (table) -> E6 isolation (table)
% -> E1 accuracy with baselines, PR curves (figure) + per-family recall.
% Each table caption names the script and JSON that produced it.
\label{sec:eval}
TODO

\section{Discussion}
% Honest items: RF beats the MLP on the same features (what that means for a
% sandboxed detector: the sandbox is model-agnostic; tree models fit too);
% per-event online alerting vs per-trace scoring; STIDE training; the openat
% volume and kernel-side pre-filtering as the lever behind the CPU gap to
% Falco/Tetragon; Tetragon's security_file_open semantics; IPv6; wasmtime
% version; single dataset (ADFA-LD, i386 numbers) if E8 is cut.
\label{sec:discussion}
TODO

\section{Related Work}
\label{sec:related}
% Prior-art check done 2026-10-25 (web search; asta flaky). No priority
% wording: the nearest work is listed and the difference stated.

\paragraph{Syscall-based host intrusion detection.}
Modelling a process by its system-call stream goes back to STIDE's fixed-length
n-gram databases~\cite{stide} and the pH delay-based response that followed
it~\cite{ph}. ADFA-LD~\cite{adfa-ld} and the semantic pattern models built on
it~\cite{creech} remain the common benchmark for this line, and sequence
models from LSTM language models~\cite{kim-lstm} and DeepLog~\cite{deeplog}
to provenance-graph detectors such as StreamSpot~\cite{streamspot} and
Unicorn~\cite{unicorn} have raised the ceiling; LID-DS~\cite{lidds} offers a
more recent recording. \sys contributes no new detector. It takes the standard
dataset and the standard baselines and reports what the \emph{deployed}
module scores, which, as \S\ref{sec:eval:e1} shows, is below a random forest
on the same features. That is the point: the number is the module's own.

\paragraph{eBPF runtime security monitors.}
Falco~\cite{falco}, Tetragon~\cite{tetragon} and Tracee~\cite{tracee} collect
process, file and network events with eBPF and evaluate rules against them;
bpftrace and the BPF tool suite~\cite{bpftrace,ebpf-intro} do the same for
ad-hoc tracing, and commercial agents such as Sysdig and
Pixie~\cite{sysdig,pixie} package it for fleets. In all of them the rule logic
is trusted: Falco rules are a filter language interpreted by the agent and
its plugins are native shared objects loaded into the agent's address space;
Tetragon policies are compiled into eBPF programs and a user-space agent that
run with the agent's privileges. The eBPF verifier is itself a sandbox, but
one with a history of bypasses~\cite{ebpf-bugs}, and it constrains what a rule
may compute far more than it constrains what a rule may see. \sys keeps the
eBPF side minimal (fixed hooks, a fixed record) and moves everything a rule
author writes into Wasm. We measure against Falco and Tetragon directly
(\S\ref{sec:eval:e4}, \S\ref{sec:eval:e5}) rather than quoting their published
figures.

\paragraph{WebAssembly as an isolation boundary.}
WebAssembly's linear memory and import-only interface~\cite{wasm-spec,wasi}
have made it a software-fault-isolation substrate well beyond the browser:
RLBox sandboxes third-party libraries in Firefox~\cite{rlbox},
Kolosick et al.\ drive the transition cost toward zero~\cite{wasmboxc},
ERIM reaches similar goals with protection keys instead~\cite{erim},
Sledge~\cite{sledge} and Lucet~\cite{lucet} run it at the edge, and
Proxy-Wasm~\cite{proxy-wasm} is the production precedent for untrusted
plugins inside a network proxy. Wasmtime~\cite{wasmtime} supplies the fuel,
epoch and memory limits \sys relies on; CT-Wasm~\cite{ct-wasm} addresses a
side-channel threat we leave out of scope. gVisor and Kata~\cite{gvisor,kata}
isolate whole workloads at a much higher cost than a per-rule instance. \sys
applies this substrate to detection rules and asks the two questions these
systems answer for their own domains: what it costs and whether it holds.

\paragraph{WebAssembly and eBPF together.}
The closest prior work combines the two technologies in the opposite
direction. Wasm-bpf~\cite{wasm-bpf} packages eBPF programs and their
user-space loaders as Wasm modules so that kernel tooling can be deployed like
cloud artifacts; Abbadini et al.~\cite{abbadini-ebpf-wasm} use eBPF to
enforce finer file-system policies on a Wasm runtime's host calls; and
Heinrich et al.~\cite{heinrich-wasm-ids} detect anomalies \emph{in} Wasm
applications from their WASI call stream. In \sys, eBPF is the sensor and Wasm
is the sandbox for the detection logic, and the evaluation is about that
boundary: the per-call cost of crossing it, the end-to-end cost of running a
host monitor behind it, and whether it contains a rule that misbehaves.

\section{Conclusion}
\label{sec:conclusion}
Host-based detection rules are code with kernel visibility, and today they
run as trusted code. \sys runs them in WebAssembly instead and this paper
measures what that costs and what it buys. The boundary itself is cheap:
\eTwoEngineOverheadUs{}~\si{\micro\second} per call empty,
\eTwoMlRawUs{}~\si{\micro\second} for a real rule on a binary-encoded event,
a sustained \eThreeAchievedBinFiftyKRf{} syscall events per second through
the whole pipeline, and an application cost under nginx, redis and
PostgreSQL in the same range as Falco's. It holds: ten adversarial modules
in eleven attempts were contained by memory, fuel and deadline limits with
their neighbours unaffected. And it makes a claim checkable that is
usually taken on faith: the detector the system runs is, to
\eTwoMaxDiffNumpyWasm{} for the MLP and \eOneWasmSklearnMaxDiff{} for the
forest, the detector that was evaluated, which is how we can report that
\sys's accuracy on ADFA-LD is an AUPRC of \eOneAuprcRfWasm{} with the forest
and \eOneAuprcMlp{} with the MLP it first shipped with, and how the better
model got into the module in a day. The remaining cost is in the rules, not
the sandbox, and we have named the two levers that move it.


\bibliographystyle{ACM-Reference-Format}
\bibliography{bib/references}

\appendix
\section{Open Science}
% Mandatory for AsiaCCS 2027 (<= half a page): anonymous artifact link
% (anonymous.4open.science) with code, rules, evaluation/results/2026-11/*.json,
% `make reproduce`, prebuilt Wasm modules and the ADFA-LD fetch script; which
% experiments need root/eBPF and which run anywhere (E1, E2, E6).
\label{sec:openscience}
TODO

\section{Additional Material}
% Per-family recall table, full E6 table, event record layout, feature set v1
% definition (verbatim from code), hyper-parameters of every baseline.
\label{sec:appendix}
TODO


\end{document}
